\subsubsection*{Rules of Replacement}

\begin{remark*}[Working vocabulary]
This topic uses formula contexts, logical equivalence, and replacement of
logical equivalents to justify algebraic rewriting inside larger formulas.
\end{remark*}

% BEGIN GENERATED ARTIFACT: def:rule-of-replacement-propositional-logic
\begin{definitionbox}{Definition (Rule of Replacement)}
\begin{definition}[Rule of Replacement]
\label{def:rule-of-replacement-propositional-logic}
A \emph{rule of replacement} is a schematic permission to replace a subformula
by a logically equivalent formula inside a formula context.
\end{definition}
\end{definitionbox}

\begin{remark*}[Standard quantified statement]
\[
\varphi\equiv\psi
\Longrightarrow
C[\varphi]\equiv C[\psi].
\]
\end{remark*}

\begin{remark*}[Definition predicate reading]
A rule of replacement licenses an equivalence-preserving rewrite inside a larger
formula.
\end{remark*}

\begin{remark*}[Interpretation]
Rules of replacement differ from rules of inference. Inference rules move from
premises to conclusions; replacement rules rewrite formulas without changing
truth behavior.
\end{remark*}

\begin{dependencies}
\begin{itemize}
  \item \hyperref[def:formula-context-propositional-logic]{Formula Context}
  \item \hyperref[def:logical-equivalence-propositional-logic]{Logical Equivalence}
  \item \hyperref[thm:replacement-of-logical-equivalents-propositional-logic]{Replacement of Logical Equivalents}
\end{itemize}
\end{dependencies}
% END GENERATED ARTIFACT: def:rule-of-replacement-propositional-logic

% BEGIN GENERATED ARTIFACT: thm:replacement-rules-preserve-logical-equivalence-propositional-logic
\begin{theorembox}{Theorem (Replacement Rules Preserve Logical Equivalence)}
\begin{theorem}[Replacement Rules Preserve Logical Equivalence]
\label{thm:replacement-rules-preserve-logical-equivalence-propositional-logic}
If a formula \(\psi\) is obtained from a formula \(\varphi\) by finitely many
applications of rules of replacement, then \(\varphi\equiv\psi\).

\noindent\hyperref[prf:replacement-rules-preserve-logical-equivalence-propositional-logic]{Go to proof.}
\end{theorem}
\end{theorembox}

\begin{remark*}[Standard quantified statement]
\[
\operatorname{ReplacementSequence}(\varphi,\psi)
\Longrightarrow
\varphi\equiv\psi.
\]
\end{remark*}

\begin{remark*}[Predicate reading]
A finite sequence of equivalence-preserving replacements preserves logical
equivalence from the first formula to the last formula.
\end{remark*}

\begin{remark*}[Interpretation]
Replacement rules justify algebraic manipulation of formulas. Each local rewrite
preserves truth under every truth assignment, so the entire finite rewrite chain
also preserves truth under every truth assignment.
\end{remark*}

\begin{remark*}[Exposition]
Common replacement rules include double negation, De Morgan replacement,
commutativity, associativity, distributivity, material implication, biconditional
replacement, and contraposition. Each is justified by a logical equivalence law.
\end{remark*}

\begin{dependencies}
\begin{itemize}
  \item \hyperref[def:rule-of-replacement-propositional-logic]{Rule of Replacement}
  \item \hyperref[thm:replacement-of-logical-equivalents-propositional-logic]{Replacement of Logical Equivalents}
  \item \hyperref[thm:logical-equivalence-equivalence-relation]{Logical Equivalence is an Equivalence Relation}
\end{itemize}
\end{dependencies}
% END GENERATED ARTIFACT: thm:replacement-rules-preserve-logical-equivalence-propositional-logic
